% GENERATED FILE. Edit canonical lessons, then run npm run generate:books.
% canonical-source-hash: fnv1a64:042950ed4d9e5b7f
% canonical-lessons: KA-C545-jamakhana, KA-C545-raggu, KA-C545-jhadisu, KA-C545-oge, KA-C545-ganji

\chapter{Laundry Day}
\label{ch:ka-laundry-day}
\hlchaptermodality{545}

\begin{chapteropening}
\textbf{By the end of this chapter:} \emph{I can say a rug, a quilt, to shake out, to wash clothes and starch.}

Complete the last lesson of chapter 545: I can say a rug, a quilt, to shake out, to wash clothes and starch.
\end{chapteropening}

\section[\texorpdfstring{\kn{ಜಮಖಾನ} (jamakhāna)}{ಜಮಖಾನ (jamakhāna)}]{\kn{ಜಮಖಾನ} (jamakhāna) — a rug; a woven floor carpet}

\label{lesson:KA-C545-jamakhana}



\begin{quote}
Before the new one: say the Kannada for a tube, then the Kannada for a wire.
\end{quote}

\subsection*{You'll want to know: \kn{ಜಮಖಾನ}}
\textbf{\kn{ಜಮಖಾನ}} — \emph{jamakhāna} — \textquotedblleft{}a rug; a woven floor carpet\textquotedblright{}.

\begin{quote}
\kn{ಹಬ್ಬಕ್ಕೆ} \kn{ನೆಲದ} \kn{ಮೇಲೆ} \kn{ಜಮಖಾನ} \kn{ಹಾಸಿದರು}.
\end{quote}

\begin{quote}
\emph{habbakke nelada mēle jamakhāna hāsidaru.}
\end{quote}

\begin{quote}
They spread a rug on the floor for the festival.
\end{quote}

Thick and woven, spread for guests to sit on.

\begin{cousinweb}
\textbf{\kn{ಜಮಖಾನ}} (\emph{jamakhāna}) is usually traced to a Persian word that reached Kannada through Hindustani; the exact source is not settled.
\end{cousinweb}

\begin{grammarlens}[title={What you've built}]
One more word for describing what you lived and felt.
\end{grammarlens}

\subsection*{Your turn}
\begin{itemize}
\raggedright
  \item \emph{Say it:} \emph{jamakhāna}
  \item \emph{Say it:} \emph{jamakhāna}, once more
  \item \emph{Say it:} \emph{habbakke nelada mēle jamakhāna hāsidaru.} — They spread a rug on the floor for the festival.
  \item \emph{Recall:} say the Kannada for a tube, then the Kannada for a wire, then say \emph{jamakhāna} again
\end{itemize}

\begin{tcolorbox}[breakable,colback=teal!4,colframe=teal!35!black,title={Before you move on}]
What does \kn{ಜಮಖಾನ} mean? (\textquotedblleft{}A rug; a woven floor carpet\textquotedblright{}.) Say it once more.
\end{tcolorbox}
\section[\texorpdfstring{\kn{ರಗ್ಗು} (raggu)}{ರಗ್ಗು (raggu)}]{\kn{ರಗ್ಗು} (raggu) — a quilt; a thick blanket}

\label{lesson:KA-C545-raggu}



\begin{quote}
Before the new one: say the Kannada for a wire, then the Kannada for a rug.
\end{quote}

\subsection*{You'll want to know: \kn{ರಗ್ಗು}}
\textbf{\kn{ರಗ್ಗು}} — \emph{raggu} — \textquotedblleft{}a quilt; a thick blanket\textquotedblright{}.

\begin{quote}
\kn{ಚಳಿಯಲ್ಲಿ} \kn{ಮಕ್ಕಳು} \kn{ರಗ್ಗು} \kn{ಹೊದ್ದು} \kn{ಮಲಗಿದರು}.
\end{quote}

\begin{quote}
\emph{caḷiyalli makkaḷu raggu hoddu malagidaru.}
\end{quote}

\begin{quote}
In the cold the children slept under a quilt.
\end{quote}

Heavier than \kn{ಕಂಬಳಿ}.

\begin{cousinweb}
\textbf{\kn{ರಗ್ಗು}} (\emph{raggu}) is a loan from English \emph{rug}, in the older sense of a thick travelling blanket.
\end{cousinweb}

\begin{grammarlens}[title={What you've built}]
One more word for describing what you lived and felt.
\end{grammarlens}

\subsection*{Your turn}
\begin{itemize}
\raggedright
  \item \emph{Say it:} \emph{raggu}
  \item \emph{Say it:} \emph{raggu}, once more
  \item \emph{Say it:} \emph{caḷiyalli makkaḷu raggu hoddu malagidaru.} — In the cold the children slept under a quilt.
  \item \emph{Recall:} say the Kannada for a wire, then the Kannada for a rug, then say \emph{raggu} again
\end{itemize}

\begin{tcolorbox}[breakable,colback=teal!4,colframe=teal!35!black,title={Before you move on}]
What does \kn{ರಗ್ಗು} mean? (\textquotedblleft{}A quilt; a thick blanket\textquotedblright{}.) Say it once more.
\end{tcolorbox}
\section[\texorpdfstring{\kn{ಝಾಡಿಸು} (jhāḍisu)}{ಝಾಡಿಸು (jhāḍisu)}]{\kn{ಝಾಡಿಸು} (jhāḍisu) — to shake out; to dust off}

\label{lesson:KA-C545-jhadisu}



\begin{quote}
Before the new one: say the Kannada for a rug, then the Kannada for a quilt.
\end{quote}

\subsection*{You'll want to know: \kn{ಝಾಡಿಸು}}
\textbf{\kn{ಝಾಡಿಸು}} — \emph{jhāḍisu} — \textquotedblleft{}to shake out; to dust off\textquotedblright{}.

\begin{quote}
\kn{ಅಜ್ಜಿ} \kn{ಜಮಖಾನವನ್ನು} \kn{ಹೊರಗೆ} \kn{ಝಾಡಿಸಿದರು}.
\end{quote}

\begin{quote}
\emph{ajji jamakhānavannu horage jhāḍisidaru.}
\end{quote}

\begin{quote}
Grandmother shook out the rug outside.
\end{quote}

Past (I …): \kn{ಝಾಡಿಸಿದೆ}, \emph{jhāḍiside}.

To shake or beat the dust out of something.

\begin{cousinweb}
\textbf{\kn{ಝಾಡಿಸು}} (\emph{jhāḍisu}) is usually traced to Hindustani \emph{jhāṛnā}, \textquotedblleft{}to dust, to shake off\textquotedblright{}, with the Kannada verb ending -\emph{isu}.
\end{cousinweb}

\begin{grammarlens}[title={What you've built}]
One more word for describing what you lived and felt.
\end{grammarlens}

\subsection*{Your turn}
\begin{itemize}
\raggedright
  \item \emph{Say it:} \emph{jhāḍisu}
  \item \emph{Say it:} \emph{jhāḍisu}, once more
  \item \emph{Say it:} \emph{ajji jamakhānavannu horage jhāḍisidaru.} — Grandmother shook out the rug outside.
  \item \emph{Recall:} say the Kannada for a rug, then the Kannada for a quilt, then say \emph{jhāḍisu} again
\end{itemize}

\begin{tcolorbox}[breakable,colback=teal!4,colframe=teal!35!black,title={Before you move on}]
What does \kn{ಝಾಡಿಸು} mean? (\textquotedblleft{}To shake out; to dust off\textquotedblright{}.) Say it once more.
\end{tcolorbox}
\section[\texorpdfstring{\kn{ಒಗೆ} (oge)}{ಒಗೆ (oge)}]{\kn{ಒಗೆ} (oge) — to wash clothes}

\label{lesson:KA-C545-oge}



\begin{quote}
Before the new one: say the Kannada for a quilt, then the Kannada for to shake out.
\end{quote}

\subsection*{You'll want to know: \kn{ಒಗೆ}}
\textbf{\kn{ಒಗೆ}} — \emph{oge} — \textquotedblleft{}to wash clothes\textquotedblright{}.

\begin{quote}
\kn{ಭಾನುವಾರ} \kn{ಅಮ್ಮ} \kn{ಎಲ್ಲಾ} \kn{ಬಟ್ಟೆ} \kn{ಒಗೆಯುತ್ತಾರೆ}.
\end{quote}

\begin{quote}
\emph{bhānuvāra amma ellā baṭṭe ogeyuttāre.}
\end{quote}

\begin{quote}
On Sunday mother washes all the clothes.
\end{quote}

Past (I …): \kn{ಒಗೆದೆ}, \emph{ogede}.

For clothes; dishes and hands take \kn{ತೊಳೆ}.

\begin{cousinweb}
\textbf{\kn{ಒಗೆ}} (\emph{oge}) is native Dravidian: DEDR 933 lists Kannada \emph{oge}, \textquotedblleft{}to wash\textquotedblright{}.
\end{cousinweb}

\begin{grammarlens}[title={What you've built}]
One more word for describing what you lived and felt.
\end{grammarlens}

\subsection*{Your turn}
\begin{itemize}
\raggedright
  \item \emph{Say it:} \emph{oge}
  \item \emph{Say it:} \emph{oge}, once more
  \item \emph{Say it:} \emph{bhānuvāra amma ellā baṭṭe ogeyuttāre.} — On Sunday mother washes all the clothes.
  \item \emph{Recall:} say the Kannada for a quilt, then the Kannada for to shake out, then say \emph{oge} again
\end{itemize}

\begin{tcolorbox}[breakable,colback=teal!4,colframe=teal!35!black,title={Before you move on}]
What does \kn{ಒಗೆ} mean? (\textquotedblleft{}To wash clothes\textquotedblright{}.) Say it once more.
\end{tcolorbox}
\section[\texorpdfstring{\kn{ಗಂಜಿ} (gañji)}{ಗಂಜಿ (gañji)}]{\kn{ಗಂಜಿ} (gañji) — starch; rice gruel}

\label{lesson:KA-C545-ganji}



\begin{quote}
Before the new one: say the Kannada for to shake out, then the Kannada for to wash clothes.
\end{quote}

\subsection*{You'll want to know: \kn{ಗಂಜಿ}}
\textbf{\kn{ಗಂಜಿ}} — \emph{gañji} — \textquotedblleft{}starch; rice gruel\textquotedblright{}.

\begin{quote}
\kn{ಅಮ್ಮ} \kn{ಹತ್ತಿ} \kn{ಸೀರೆಗೆ} \kn{ಗಂಜಿ} \kn{ಹಾಕಿದರು}.
\end{quote}

\begin{quote}
\emph{amma hatti sīrege gañji hākidaru.}
\end{quote}

\begin{quote}
Mother starched the cotton sari.
\end{quote}

Also the water rice is cooked in, a light food for the sick.

\begin{cousinweb}
\textbf{\kn{ಗಂಜಿ}} (\emph{gañji}) is native Dravidian: DEDR 1107 lists Kannada \emph{gañji}, \textquotedblleft{}rice-gruel, starch\textquotedblright{}.
\end{cousinweb}

\begin{grammarlens}[title={What you've built}]
One more word for describing what you lived and felt.
\end{grammarlens}

\subsection*{Your turn}
\begin{itemize}
\raggedright
  \item \emph{Say it:} \emph{gañji}
  \item \emph{Say it:} \emph{gañji}, once more
  \item \emph{Say it:} \emph{amma hatti sīrege gañji hākidaru.} — Mother starched the cotton sari.
  \item \emph{Recall:} say the Kannada for to shake out, then the Kannada for to wash clothes, then say \emph{gañji} again
\end{itemize}

\begin{tcolorbox}[breakable,colback=teal!4,colframe=teal!35!black,title={Before you move on}]
What does \kn{ಗಂಜಿ} mean? (\textquotedblleft{}Starch; rice gruel\textquotedblright{}.) Say it once more.
\end{tcolorbox}
